\section{Reuse geometry and two-layer Leech lifting}\label{sec:leech}

The endpoint and four-triangle constructions below isolate the geometry
of direction reuse: the latitude, the tail partition, and the capacity
of a prescribed label system. They give the auxiliary configurations
of sizes 197058 and 200540. Section~\ref{sec:leech-two-layer} then
uses a second layer to obtain the current 201567-point construction.

We use the integer Leech shell $\Sigma\subset\Z^{24}$ in the Golay
coordinate convention of Appendix~\ref{app:leech}. Its relevant
properties are
\begin{equation}\label{eq:leech-shell}
 |\Sigma|=196560,\qquad r\cdot r=32,\qquad
 r\cdot r'\le16\quad(r\ne r',\ r,r'\in\Sigma).
\end{equation}
Consequently $C=\Sigma/\sqrt{32}$ is a kissing configuration. A block
$S\subset\Sigma$ will be called \emph{compatible} if distinct rows
satisfy $r\cdot r'\le8$. After normalization its maximum inner product
is at most $1/4$.

\begin{lemma}[Leech block witnesses]\label{lem:leech-blocks}
There is a compatible set $S_*\subset\Sigma$ of size $496$ and there
are four pairwise disjoint compatible sets $S_0,S_1,S_2,S_3\subset\Sigma$,
each of size $496$.
\end{lemma}
\begin{proof}[Finite verification]
The five complete integer arrays are specified in
Table~\ref{tab:certificates}. They are taken from the
PackingStar coordinate data~\cite{packingstar-data}.
For each array, membership is checked against the three shell families
in Appendix~\ref{app:leech}. Each array has $496$ distinct rows, each
row has squared norm $32$, and the maximum of its $122760$ unordered
distinct-row inner products is eight. The latter four arrays have a
union of cardinality $1984$. These checks, implemented separately by
the two programs described in Section~\ref{sec:certificates}, prove
compatibility and disjointness.
\end{proof}

\subsection{The endpoint construction in dimension 25}
The normalized block $S_*/\sqrt{32}$ satisfies the $1/3$ condition of
Theorem~\ref{thm:block-lift}, since it satisfies the stronger bound $1/4$.
The resulting configuration is
\begin{align}
 X_{25}={}&\{(r/\sqrt{32},0):r\in\Sigma\setminus S_*\}\nonumber\\
 &\quad\cup\{(\sqrt3\,r/\sqrt{128},\eps/2):
                   r\in S_*,\ \eps\in\{\pm1\}\}
       \cup\{(0,1),(0,-1)\}.\label{eq:d25}
\end{align}
Its cardinality is
\[
 196064+992+2=197058,
\]
giving the historical bound $K(25)\ge197058$.
The endpoint choice permits the two poles without requiring a larger
Leech block. Distinct same-latitude directions have inner product at
most $7/16$, whereas the opposite-latitude copies of one direction
have inner product exactly $1/2$. This distinction will also determine
the local geometry of $X_{25}$ in Section~\ref{sec:deformations}.

\subsection{The tail geometry in dimension 27}
Let
\begin{equation}\label{eq:D3}
 V=\{\pm e_i\pm e_j:1\le i<j\le3\},\qquad
 R=\begin{pmatrix}
  1/\sqrt2&-1/\sqrt2&0\\
  1/\sqrt2&1/\sqrt2&0\\
  0&0&1
 \end{pmatrix}.
\end{equation}
The twelve roots in $V$ form a cuboctahedron. Put
\begin{equation}\label{eq:tailsets}
 T=V/\sqrt6,\qquad P=RV/\sqrt2.
\end{equation}
Thus $T$ supplies mixed tails of squared norm $1/3$, and $P$ supplies
twelve unit pure tails. The relative rotation is through $45$ degrees
in the first two coordinates.

\begin{lemma}\label{lem:tails}
Distinct members of $T$ have inner products in
$\{-1/3,-1/6,0,1/6\}$, and $P$ is a kissing configuration. Moreover,
\begin{equation}\label{eq:tail-cross}
 \max_{p\in P,\,t\in T}\ip{p}{t}
  =\frac{1+1/\sqrt2}{\sqrt{12}}<\frac12.
\end{equation}
The roots admit a partition into four zero-sum triples, each of which
has pairwise inner product $-1$.
\end{lemma}
\begin{proof}
Every root has squared norm two. Distinct roots have inner products
$-2,-1,0,$ or $1$, giving the first assertions by scaling and
orthogonality of $R$. The absolute coordinates of $Rv$, up to order,
are either $(\sqrt2,0,0)$ or $(1,1/\sqrt2,1/\sqrt2)$. Its largest inner
product with a root is obtained by choosing two coordinates with the
appropriate signs. The maximum is $1+1/\sqrt2$, giving
\eqref{eq:tail-cross}. The strict inequality follows, for example, from
$18+12\sqrt2<36$ after squaring the equivalent inequality
$\sqrt6+2\sqrt3<6$.

An explicit partition is given by the rows of
\begin{equation}\label{eq:triangles}
\begin{array}{rrr}
 (-1,0,-1)&(1,-1,0)&(0,1,1)\\
 (0,-1,-1)&(-1,0,1)&(1,1,0)\\
 (-1,-1,0)&(0,1,-1)&(1,0,1)\\
 (-1,1,0)&(1,0,-1)&(0,-1,1).
\end{array}
\end{equation}
The roots in a row sum to zero. Their equal squared norms then force
all three mutual inner products to be $-1$.
\end{proof}

Write $T_b$ for row $b$ of~\eqref{eq:triangles}, divided by $\sqrt6$.
Thus $T$ is the disjoint union of $T_0,T_1,T_2,T_3$. A block assigned
to $T_b$ is used at three labels, while its original equatorial copy
is removed once. This is the multiplicity-three counterpart of the
antipodal lift.

\begin{proposition}[Four-triangle lift]\label{prop:triangles}
For pairwise disjoint compatible subsets $S_0,S_1,S_2,S_3$ of $\Sigma$,
the following three sets form a kissing configuration:
\begin{align}
 Z&=\{(r/\sqrt{32},0):r\in\Sigma\setminus\textstyle\bigcup_bS_b\},
 \nonumber\\
 M&=\bigcup_{b=0}^{3}\{(r/\sqrt{48},t):r\in S_b,\ t\in T_b\},
 \label{eq:d27}\\
 P_0&=\{(0,p):p\in P\}.\nonumber
\end{align}
Its cardinality is $196560+2\sum_b|S_b|+12$.
\end{proposition}
\begin{proof}
Use Theorem~\ref{lem:lift} with $h^2=1/3$ and $a=\sqrt{2/3}$.
The mixed vectors have norm one since $32/48+1/3=1$.
For mixed points, multiply their inner-product condition by $48$:
\begin{equation}\label{eq:raw-mixed}
 r\cdot r'+48\ip{t}{t'}\le24.
\end{equation}
If $t=t'$, the raw vectors are distinct and lie in one compatible
block, so the left side is at most $8+16=24$. If $r=r'$ and $t\ne t'$,
disjointness of the blocks places both tails in one triangle, giving
$32-8=24$. If both directions and labels differ, the universal shell
bound and Lemma~\ref{lem:tails} give $16+8=24$ as an upper bound.
These three cases cover every distinct mixed pair. Pure--mixed
compatibility follows from~\eqref{eq:tail-cross}; all remaining pairs
are covered by Theorem~\ref{lem:lift}.

For $B=\sum_b|S_b|$, the layer sizes are $196560-B$, $3B$, and $12$.
They are disjoint, and their parametrizations are injective, as in
\eqref{eq:lift-count}. Their sum is the asserted count.
\end{proof}

Applying Proposition~\ref{prop:triangles} to Lemma~\ref{lem:leech-blocks}
gives
\[
 |Z|=194576,\qquad |M|=5952,\qquad |P_0|=12,
\]
and hence $|X_{27}|=200540$. The comparison allocation uses block
multiplicities $(3,3,2,2,2)$, whereas~\eqref{eq:triangles} uses
$(3,3,3,3)$. Both occupy twelve labels. The latter allocation uses
$496$ fewer original directions and therefore preserves $496$ more
equatorial points.

\subsection{Capacity and its equality case}
The preceding count has an exact interpretation even when the label
sets are not arranged in disjoint blocks. Let $A_i\subseteq\Sigma$
be the rows assigned to tail $t_i\in T$. Assume that the resulting
lifting is valid, retains $\Sigma\setminus\bigcup_iA_i$ at the equator,
and contains all twelve pure tails.

\begin{proposition}\label{prop:capacity}
If $|A_i|\le\alpha$ for each of the twelve labels, then
\begin{equation}\label{eq:capacity}
 |X|\le196560+8\alpha+12.
\end{equation}
Equality holds precisely when every used row occurs at three labels
and every label contains $\alpha$ rows. For $\alpha=496$, the
four-triangle construction attains equality.
\end{proposition}
\begin{proof}
Proposition~\ref{prop:multiplicity} gives maximum multiplicity three.
Writing $m(r)$ for the multiplicity of a used row,
\[
 |X|-196560-12
 =\sum_r(m(r)-1)
 \le\frac23\sum_rm(r)
 =\frac23\sum_i|A_i|\le8\alpha.
\]
The first inequality is an equality exactly when all $m(r)=3$;
the last is an equality exactly when all labels are full. The explicit
assignment has both properties.
\end{proof}

More precisely, let $n_j$ count the used rows with multiplicity $j$,
and let $c_i=\alpha-|A_i|$. The deficit is
\begin{equation}\label{eq:deficit}
 196560+8\alpha+12-|X|
 =\frac23\sum_i c_i+\frac23n_1+\frac13n_2.
\end{equation}
For the comparison allocation, $n_2=3\cdot496$ and all $c_i=0$, so
the deficit is exactly $496$. The capacity assumption in
Proposition~\ref{prop:capacity} concerns the prescribed label system;
the block witnesses supply $\alpha=496$ without requiring a bound on
the largest possible compatible subset of the Leech shell.

\subsection{A two-layer configuration with 201567 points}
\label{sec:leech-two-layer}
The larger construction retains the first-layer heads of the public
201566-point contribution~\cite{lindow27} and replaces its second
layer by a 311-owner witness. We give the pair conditions explicitly.
For this subsection use squared norm 32 and threshold 16 throughout.
Write $\Sigma_{27}$ for the Leech shell generated in the Golay
coordinate basis of the two-layer data. All heads and owners below
use that same basis. The shell argument in Appendix~\ref{app:leech}
applies to its verified Golay generator.

The first-layer data are 2258 distinct integer heads $Y$ of squared
norm 192, each assigned to one of four triples. In the order of the
stored labels, these triples of roots are
\[
\begin{array}{c|rrr}
0&(1,1,0)&(-1,0,-1)&(0,-1,1)\\
1&(1,-1,0)&(-1,0,1)&(0,1,-1)\\
2&(-1,-1,0)&(1,0,-1)&(0,1,1)\\
3&(-1,1,0)&(1,0,1)&(0,-1,-1).
\end{array}
\]
For a head $y$ in triple $b$, take the three points
\begin{equation}\label{eq:d27-first}
 (y/3,\,4v/\sqrt3)\qquad(v\text{ in triple }b).
\end{equation}
The pure-tail points are $(0,4Rv)$ for all twelve roots
$v\in V$, with $R$ as in~\eqref{eq:D3}.

A first-layer owner is a shell vector $z\in\Sigma_{27}$ with $y\cdot z>48$.
The finite data have 2090 heads with exactly one owner and 168
heads with none; all 2090 owners are distinct. The former are
of the form $y=3u+v$ with $u\in\Sigma_{27}$, $v^2=48$, and
$u\cdot v=-24$. The latter are twice norm-48 vectors. These
descriptions explain the candidate families; verification recovers
the owners by checking the whole shell.
For distinct first-layer heads, the stored data satisfy
\begin{equation}\label{eq:d27-first-threshold}
 y\cdot y'\le
 \begin{cases}
 48,&\text{on the same triple},\\
 96,&\text{on different triples}.
 \end{cases}
\end{equation}
Retaining all nonowners at the equator gives
$196560-2090+3(2258)+12=201256$ points.

The additional data are 311 distinct owners $u\in\Sigma_{27}$, with a
label $\ell(u)\in\{1,2,3\}$. None is a first-layer owner. Replace
the equatorial point $(u,0)$ by
\begin{equation}\label{eq:d27-second}
 (\sqrt3\,u/2,\,\eps\sqrt8\,e_{\ell(u)}),
 \qquad \eps\in\{-1,1\}.
\end{equation}
The three labels have respectively 105, 103, and 103 owners.
The checked inequalities are
\begin{align}
 y\cdot u&\le32&&\text{for every first-layer head }y,\label{eq:d27-interlayer}\\
 u\cdot u'&\le8&&\text{if }\ell(u)=\ell(u'),\quad u\ne u',\label{eq:d27-sameline}\\
 u\cdot u'&\le16&&\text{if }\ell(u)\ne\ell(u').\label{eq:d27-crossline}
\end{align}

\begin{theorem}\label{thm:d27}
The two layers just specified, their retained equator, and their
twelve pure-tail points form a kissing configuration of 201567
points in dimension 27.
\end{theorem}
\begin{proof}
All points have squared norm 32. An old equatorial point and a
first-layer cap have product at most $48/3=16$. The two
thresholds in~\eqref{eq:d27-first-threshold} give respectively
$48/9+32/3=16$ and $96/9+16/3=16$.
Two different caps of the same head have product
$192/9-16/3=16$.
The rotation estimate of Lemma~\ref{lem:tails}, scaled by 32,
checks the first layer against all pure tails.

A second-layer cap has product at most
$(\sqrt3/2)16=8\sqrt3<16$ with every retained equatorial
point. Its own owner is deleted. Against a first-layer cap, the
head product and tail product sum to at most
\[
 \frac{\sqrt3}{6}(32)+\frac{4\sqrt8}{\sqrt3}
 =\frac{16\sqrt3+8\sqrt6}{3}<16.
\]
The last inequality is equivalent to
$2\sqrt3+\sqrt6<6$ and follows by squaring positive quantities.
Second-layer caps on distinct owners of the same label have
product at most $(3/4)8+8=14$; distinct labels give at most
$(3/4)16=12$. The two caps of one owner have product
$(3/4)32-8=16$. A second-layer cap and a pure tail have
product at most $\sqrt8\sqrt{32}=16$.
All old--old pairs retain their original bounds.

The head squared norms in the equator, first layer, second layer,
and pure-tail layer are respectively $32$, $64/3$, $24$, and zero,
so different layers cannot coincide. Within a layer, distinct
heads and tail labels give distinct points. Each of the 311
new owners is removed once and inserted twice. Hence
\[
 196560-2090-311+3(2258)+2(311)+12=201567.
\]
Dividing all points by $\sqrt{32}$ proves the unit-sphere claim.
\end{proof}

The finite optimization behind the second layer has a particularly
small geometric description. Keep only shell vectors that are
not first-layer owners and satisfy~\eqref{eq:d27-interlayer}.
Connect two such vectors when their product exceeds eight.
One seeks a largest subset that can be split into three independent
sets of this graph. Assign each independent set to one coordinate
line. Distinct-label products are already controlled by the
Leech-shell bound. The published 310-owner witness and the new
311-owner witness are feasible assignments in this same model;
each extra owner contributes one net point.
